\section{Methodology}

The project will be developed in incremental and iterative stages, ensuring the methodological rigor characteristic of continuous optimization.

\subsection{Directed Study and Literature Review}
Initially, an in-depth study of the trust-region methods for nonsmooth optimization \cite{leconte2023complexity}, nonmonotone algorithmic behavior, and indefinite proximal gradient methods \cite{Leconte_2024} developed by the group at Polytechnique Montréal will be conducted. In parallel, there will be focused training on the architecture of \texttt{JuliaSmoothOptimizers} to ensure that subsequent developments are aligned with the software engineering best practices of the community.

\subsection{Development of Active Hybrid Algorithms}
The research methodology will focus on the design of a hybrid algorithmic model. The component developed in Brazil (NSPG + combinatorial searches) will act in the initial iterations of the optimization process, promoting sparsity and rapidly identifying an active set (the non-zero variables).

Additionally, we will investigate the use of structured second-order matrices in the proximal operator subproblem (Variable Metric Proximal operators). While recent approaches utilizing simple indefinite diagonal metrics have been shown to underperform compared to standard spectral methods \cite{Leconte_2024}, leveraging slightly denser structures, such as tridiagonal matrices, could provide a significant convergence boost over the spectral (identity multiple) model. In fact, our preliminary experiments on L0-regularized least squares with tridiagonal structures have already shown great promise, demonstrating that these subproblems remain computationally fast while offering high precision in support recovery.

Once the active set is identified by the global preprocessors, the algorithm will leverage the coordinate descent and combinatorial searches to strictly refine the solution over the non-zero variables. This \textit{active-set} strategy drastically reduces the computational complexity, combining the rapid landscape navigation of structured proximal operators with the precise local convergence and support isolation of the CD refiners.

\subsection{Numerical Evaluation and Benchmarking}
The implementation will be done in the Julia language, due to its high performance, which aligns with the scope of the project in Brazil and the JSO infrastructure in Canada. Extensive experiments will be conducted comparing the proposed hybrid method against state-of-the-art algorithms (such as PGCCD and the L0Learn package) on high-dimensional synthetic instances and real statistical learning datasets. The evaluation metrics will include CPU time, support precision, scalability, and the final solution quality in terms of the restricted objective function.
